% Section 8 --- Responsible Disclosure (~0.4 pages)
%
% REVIEW STATUS: First complete draft (2026-05-28).
% Vendor response table: placeholders --- to be filled at Week 14--23 (Task 9.x).
% Actual D1/D2 dates: to be filled after initial outreach (Task 3.3, Week 6).
% OTS proof hashes: to be filled after stamp_attempt.sh runs (Task 3.4).
%
% NOTE: §4.4 [DEPLOYED_APP] disclosure (D3) is not applicable in this paper's
% final form. The disclosure covers D1 (shared sequencer vendors) and D2
% (preconfirmation protocol vendors) only.

\section{Responsible Disclosure}
\label{sec:disclosure}

\subsection{Disclosure Scope}
\label{sec:disclosure-scope}

We identified two vendor categories whose deployed systems are directly
affected by the structural findings in this paper, and disclosed to each
prior to submission.

\paragraph{$D_1$ --- Shared sequencer vendors.}
The PEFO structural conditions (C1/C2/C3, \S\ref{sec:pefo-structural}) are
confirmed on three deployed shared sequencer platforms by primary source
documentation. We disclosed to:

\begin{itemize}
  \item \textbf{Astria} --- Lazy sequencer confirmed verbatim:
    \emph{``provides a guarantee on the ordering of transactions in a block,
    but it doesn't execute the STF of any given rollup''}~\cite{Astria-TxFlow}.
    The Astria README's claim of \emph{``atomic cross-rollup
    composability''}~\cite{Astria-README} directly instantiates the
    A2/A4 conflation this paper analyzes.
  \item \textbf{Espresso Systems} --- Lazy sequencer confirmed verbatim:
    \emph{``L2 validators receive blocks and execute the state transition
    functions for their rollups''} (after HotShot confirmation)~\cite{Espresso-Architecture}.
  \item \textbf{Radius} --- Sequencer model under review; disclosed as
    precautionary measure pending full architectural analysis.
\end{itemize}

The $D_1$ disclosure communicates that: (a) their platforms satisfy the
PEFO structural conditions for any cross-rollup intent-filling application
deployed on them; (b) the A2 gap cannot be closed by bond-based mechanisms
alone (Theorem~\ref{thm:a2-necessary}); and (c) we demonstrate the gap
on an Astria devnet testbed (\S\ref{sec:pefo}).

\paragraph{$D_2$ --- Preconfirmation protocol vendors.}
The PEO attack (\S\ref{sec:peo}) is demonstrated on a faithful reimplementation
of the mev-commit preconfirmation protocol~\cite{MevCommit-Docs}. We disclosed to:

\begin{itemize}
  \item \textbf{Primev (mev-commit)} --- PEO demonstrated via Foundry PoC;
    \texttt{OpenedCommitment.revertingTxHashes} field confirms inclusion-only
    commitment; $\Delta\Pi = \max(\delta - s \cdot 1.05, 0)$ payoff model
    with concrete $B = 1$ ETH, $s = 0.05$ ETH, $\delta = 0.12$ ETH worked example.
  \item \textbf{Chainbound (Bolt)} --- Bolt was archived in February 2026
    (v0.4.0-alpha). We notified Chainbound as a precautionary measure,
    noting that the PEO analysis applies to any inclusion-based bonded
    preconfirmation system.
\end{itemize}

\subsection{Disclosure Rationale}
\label{sec:disclosure-rationale}

We follow responsible disclosure norms for system security
research~\cite{iso29147}:

\begin{enumerate}
  \item \textbf{Structural, not exploit-specific.} The PEO and PEFO
    findings are structural properties of the commitment model, not
    vulnerabilities in specific contract implementations. The disclosures
    communicate the structural condition so vendors can assess whether
    their applications are affected.
  \item \textbf{Pre-publication notification.} All disclosures were sent
    at least 90 days before the submission deadline, providing vendors
    opportunity to respond, patch, or coordinate public statements.
  \item \textbf{Testbed-only PoC.} All proofs-of-concept use devnet
    infrastructure with no real value at risk. No mainnet funds were accessed,
    transferred, or put at risk at any point.
  \item \textbf{No weaponization.} The PoC scripts in Appendix~\ref{app:poc}
    require local deployment of test contracts. They are not usable against
    live systems without significant adaptation, and we have not provided
    such adaptation.
\end{enumerate}

\subsection{Vendor Responses}
\label{sec:disclosure-responses}

Table~\ref{tab:vendor-responses} summarizes the disclosure timeline and
vendor responses as of the submission date.

\begin{table}[h]
\centering
\small
\caption{Disclosure timeline and vendor responses. $D_1$/$D_2$ = date of
  initial outreach; ``Resp.'' = date of substantive response (if any);
  ``Outcome'' = vendor's characterization of the finding.
  OTS proofs in \texttt{disclosure/attempts/}.}
\label{tab:vendor-responses}
\begin{tabular}{@{}lllll@{}}
\toprule
Vendor & Category & $D_x$ & Resp. & Outcome \\
\midrule
Astria       & $D_1$ & [date] & [date] & pending \\
Espresso     & $D_1$ & [date] & [date] & pending \\
Radius       & $D_1$ & [date] & --- & precautionary \\
Primev       & $D_2$ & [date] & [date] & pending \\
Chainbound   & $D_2$ & [date] & --- & archived (Feb 2026) \\
\bottomrule
\end{tabular}
\end{table}

\noindent Vendor responses are reproduced (with vendor permission) in
Appendix~\ref{app:disclosure}. All outreach attempts are timestamped
per the scheme described in \S\ref{sec:disclosure-timestamps}.

\subsection{Timestamping and Verifiability}
\label{sec:disclosure-timestamps}

All disclosure attempts are timestamped using a two-layer scheme to
provide independent verifiable evidence of timing:

\begin{enumerate}
  \item \textbf{Per-attempt OpenTimestamps proof.} Each outreach message
    is hashed (SHA-256) and submitted to the OpenTimestamps calendar
    network~\cite{opentimestamps} immediately after sending. The resulting
    \texttt{.ots} proof file anchors the content hash to a Bitcoin block,
    providing a commitment that the message existed at or before the block
    timestamp. Proof files are stored in
    \texttt{disclosure/attempts/<date>\_<vendor>\_<channel>/content.txt.ots}.

  \item \textbf{Git commit anchor.} Each outreach is committed to the
    public artifact repository~\cite{artifact-repo} immediately after
    sending, creating a second independent timestamp via the GitHub
    commit timestamp. The commit hash for each disclosure attempt is
    recorded in \texttt{disclosure/disclosure\_commitments.json}.
\end{enumerate}

\noindent The combination of (1) and (2) allows independent verification that
the disclosure content was fixed before the submission date: the OTS proof
is anchored in Bitcoin (tamper-evident, cannot be backdated), and the git
commit history is publicly auditable.

\paragraph{Verification.}
\begin{verbatim}
# Verify a specific disclosure attempt:
ots verify disclosure/attempts/<attempt>/content.txt.ots
# -> Got 1 attestation(s) from calendar...
# -> Bitcoin block <N> attests data existed before <timestamp>

# Check SHA-256 integrity:
shasum -a 256 disclosure/attempts/<attempt>/content.txt
# -> matches content.txt.ots proof
\end{verbatim}

The aggregated disclosure record is in
\texttt{disclosure/disclosure\_commitments.json}, containing per-attempt
entries with fields: \texttt{vendor}, \texttt{date}, \texttt{channel},
\texttt{git\_commit\_hash}, \texttt{opentimestamps\_proof},
\texttt{content\_sha256}.

\subsection{Fallback Policy}
\label{sec:disclosure-fallback}

This paper will be published regardless of vendor response. Our disclosure
is a notification, not a permission request. If a vendor provides a
substantive technical response before submission, we incorporate it in
\S\ref{sec:disclosure-responses} with their permission. If a vendor
requests additional time beyond the 90-day window, we evaluate the request
on a case-by-case basis.

The structural findings --- Theorems~\ref{thm:a2-necessary}
and~\ref{thm:bond-insufficient} --- are mathematical results that are not
subject to vendor dispute on empirical grounds; they hold for any system
satisfying the definitional preconditions. Vendor responses may provide
context about operational mitigations, design rationale, or planned changes,
which we would document accurately in the vendor response table.
